\clearpage
\VerdierPhysicalPageAnchor{0066}%
\VerdierSemanticAnchor{ch1-2}
\section{Complexes in an Additive Category}

\VerdierSemanticAnchor{ch1-2-1}
\subsection{\texorpdfstring{$n$}{n}-Fold Complexes}

\VerdierSemanticAnchor{ch1-2-1-1}
\noindent\textbf{2.1.1.}\quad
Let us first fix some notation. We denote by $[n]$ the set

\VerdierSemanticAnchor{ch1-2-1-1-set}
\begin{equation*}
  [n]=\{p\in\mathbb Z\mid 1\leq p\leq n\}.
\end{equation*}

\VerdierSemanticAnchor{ch1-2-1-1-free-group}
We denote by $\mathbb Z[n]$ the free abelian group generated by the set
$[n]$, and for every $i\in[n]$ we denote by $e_i$ the generator of
$\mathbb Z[n]$ corresponding to $i$. Let
$\varphi\colon[n]\longrightarrow[m]$ be a map; we also denote by
$\varphi\colon\mathbb Z[n]\longrightarrow\mathbb Z[m]$ the corresponding
group homomorphism. In this paragraph, $C$ denotes an additive category.

\VerdierSemanticAnchor{ch1-2-1-2}
\VerdierSemanticAnchor{ch1-2-1-2-definition}
\noindent\textbf{Definition 2.1.2.}\quad
An \emph{$n$-fold complex of objects of $C$} is an object consisting of:

\begin{enumerate}
\renewcommand{\labelenumi}{\arabic{enumi})}
\item \VerdierSemanticAnchor{ch1-2-1-2-objects}%
  for every $\sigma\in\mathbb Z[n]$, an object $X^\sigma$ of $C$;
\item \VerdierSemanticAnchor{ch1-2-1-2-differentials}%
  for every $\sigma\in\mathbb Z[n]$ and every $i\in[n]$, a morphism
  \VerdierSemanticAnchor{ch1-2-1-2-differential-map}%
  \begin{equation*}
    d^{\sigma,i}\colon X^\sigma\longrightarrow X^{\sigma+e_i},
  \end{equation*}
  these morphisms being subject to the following conditions:
\end{enumerate}

\VerdierSemanticAnchor{ch1-2-1-2-eq1}
\begin{equation}
  \begin{aligned}
    d^{\sigma+e_i,i}d^{\sigma,i}&=0
      &&\forall\sigma\in\mathbb Z[n],\ \forall i\in[n],\\
    d^{\sigma+e_j,i}d^{\sigma,j}
      &=d^{\sigma+e_i,j}d^{\sigma,i}
      &&\forall\sigma\in\mathbb Z[n],\ \forall i,j\in[n].
  \end{aligned}
  \tag{2.1.2.1}
\end{equation}

\VerdierSemanticAnchor{ch1-2-1-3}
\noindent\textbf{2.1.3.}\quad
An $n$-fold complex of objects of $C$ will usually be called an
\emph{$n$-fold complex in $C$}. Let $(X^\sigma,d^{\sigma,i})$ be an
$n$-fold complex in $C$.
\VerdierSemanticAnchor{ch1-2-1-3-component}
The object $X^\sigma$ is called the \emph{component of degree $\sigma$}
of the complex $(X^\sigma,d^{\sigma,i})$.
\VerdierSemanticAnchor{ch1-2-1-3-differential}
The morphism $d^{\sigma,i}$ is called the \emph{differential of degree
$\sigma$ in direction $i$} of the complex $(X^\sigma,d^{\sigma,i})$.
\VerdierSemanticAnchor{ch1-2-1-3-notation}
Let $Y$ be an $n$-fold complex in $C$. The component of degree $\sigma$
of $Y$ is denoted $Y^\sigma$. The differential of degree $\sigma$ in
direction $i$ of $Y$ is denoted $d_Y^{\sigma,i}$.

\VerdierSemanticAnchor{ch1-2-1-4}
\VerdierSemanticAnchor{ch1-2-1-4-definition}
\noindent\textbf{Definition 2.1.4.}\quad
Let $Y$ and $Y'$ be two $n$-fold complexes in $C$. A \emph{morphism of
complexes from $Y$ to $Y'$} is a family of morphisms in $C$

\VerdierSemanticAnchor{ch1-2-1-4-family}
\begin{equation*}
  \bigl(f^\sigma\colon Y^\sigma\longrightarrow Y'^{\sigma}\bigr)_{
    \sigma\in\mathbb Z[n]}
\end{equation*}
that commute with the differentials, i.e. such that

\VerdierSemanticAnchor{ch1-2-1-4-eq1}
\begin{equation}
  f^{\sigma+e_i}d_Y^{\sigma,i}
  =d_{Y'}^{\sigma,i}f^\sigma.
  \tag{2.1.4.1}
\end{equation}
